\section{Distributed point anchors in dimensions three and four}\label{sec:points}
We prove the two smaller full counts in Theorem~\ref{thm:checker}.
The finite point estimate uses all constant modes and counting factors.
Its dimension dependence is classical in character; the exact inequalities
needed here are derived rather than imported from a continuum estimate.

\begin{lemma}[Finite balanced-box point estimates]\label{lem:point}
Let $p=2$ or $3$, let $L\geq1$ be an integer, and let $Q$ be a product of
$p$ integer paths with lengths $\ell_i\in[L/2,L]$. Use counting norm and
internal edge energy $\En_Q$, each path edge once. At any $b\in Q$,
$f(b)=0$ implies
\begin{equation}\label{eq:pointbox}
 \norm f_Q^2\leq
 \begin{cases}
 13L^2(1+\ln L)\En_Q(f),&p=2,\\
 113L^3\En_Q(f),&p=3.
 \end{cases}
\end{equation}
Here $\norm f_Q^2$ denotes $\sum_{y\in Q}|f(y)|^2$.
\end{lemma}
\begin{proof}
Put $N=\prod_i\ell_i$. If $N=1$ the anchored function is zero, so assume
$N>1$. In particular $L\geq2$. The tensor basis from
Lemma~\ref{lem:neumann} has point values and eigenvalues
\[
 |\phi_\nu(b)|^2\leq2^p/N,\qquad
 \lambda_\nu=4\sum_i\sin^2\frac{\pi\nu_i}{2\ell_i}
       \geq\frac{4|\nu|^2}{L^2}\quad(\nu\ne0),
 \qquad0\leq\nu_i<\ell_i.
\]
A singleton factor forces its index to zero. All indices are nonnegative
cosine indices; no signed-frequency multiplicity is suppressed.
On the shell $r=\max_i\nu_i\geq1$, for $p=2$ there are at most
$2r+1\leq3r$ indices, and for $p=3$ at most
$3r^2+3r+1\leq7r^2$. Since $|\nu|^2\geq r^2$ and $r\leq L-1$,
\[
 \sum_{\nu\ne0}|\nu|^{-2}\leq
 \begin{cases}3(1+\ln L),&p=2,\\7L,&p=3.\end{cases}
\]
The first bound follows from
$\sum_{r=1}^{L-1}r^{-1}\leq1+\ln L$ by integration; the second is the
finite sum of at most $L-1$ terms each bounded by seven.
Consequently the nonconstant diagonal inverse satisfies
\begin{equation}\label{eq:finiteGreen}
 G_Q(b,b):=\sum_{\nu\ne0}\frac{|\phi_\nu(b)|^2}{\lambda_\nu}
 \leq\begin{cases}
 3L^2(1+\ln L)/N\leq12(1+\ln L),&p=2,\\
 14L^3/N\leq112,&p=3,
 \end{cases}
\end{equation}
using $N\geq L^p/2^p$.
Expand $f=\sum a_\nu\phi_\nu$ in this counting-orthonormal basis.
Cauchy--Schwarz gives
$|f(b)-\bar f|^2\leq G_Q(b,b)\En_Q(f)$, where $\bar f$ is its arithmetic
mean. The constant-mode mass is $N|\bar f|^2$, not $|\bar f|^2$.
If $f(b)=0$, this mass is at most $12L^2(1+\ln L)\En_Q$ when $p=2$,
and $112L^3\En_Q$ when $p=3$, because $N\leq L^p$.
The positive spectral gap separately bounds the variance by $L^2\En_Q/4$.
Adding the orthogonal constant and variance parts proves \eqref{eq:pointbox}:
$12+1/4<13$, and $112L^3+L^2/4\leq113L^3$.
The $N=1$ case includes $L=1$ and any permissible all-singleton product.
Thus no nonexistent positive gap or lost normalization is used. Point values
were bounded uniformly, so a corner anchor is permitted without an additional
boundary trace assumption.
\end{proof}

\subsection{The full strong graph in original coordinates}
Let $p=D-1$, with $D=3$ or $4$. Fix an original scalar spatial/channel
component and its strong-source parity $A$. For every residue $s$ of that
parity set, with $y\in(\Z/n\Z)^p$,
\begin{equation}\label{eq:paircoordinates}
 x_s(y)=(y_1,\ldots,y_p,s-\textstyle\sum_i y_i)\pmod n,
 \quad u_s(y)=z(x_s(y)),\quad v_s(y)=z(x_s(y)+e_D).
\end{equation}
The first map is a bijection onto sum layer $s$; the shifted map is a
bijection onto layer $s+1$. Even $n$ makes residue parity agree with original
site parity. The $n/2$ pairs $\{s,s+1\}$ cover all sum residues disjointly,
including $\{n-1,0\}$ for odd source parity. Their pair norms therefore sum
to the original scalar counting norm exactly.

The actual strong $+e_D$ edge is $u_s(y)$--$v_s(y)$. For $1\leq i\leq p$,
$x_s(y+e_i)+e_D=x_s(y)+e_i$, so the other strong edges are
$u_s(y)$--$v_s(y+e_i)$. They are \emph{all} $D$ outgoing strong labels,
not a surrogate subset with reassigned endpoint weights. Their exact energy is
\begin{equation}\label{eq:pairenergy}
 E_s=\sum_y|u_s(y)-v_s(y)|^2+
          \sum_{y,i=1}^p|u_s(y)-v_s(y+e_i)|^2.
\end{equation}
A cross edge followed by a rung connects $u_s(y)$ to $u_s(y+e_i)$; all source
translations are generated and every target has a rung, proving connectivity.
At $n=2$ the heads $y,y+e_i$ are distinct. The induced source-torus energy
may repeat an undirected neighbor pair, but every original forward label is
retained. Translations remain bijections.

Insert $v_s(y+e_i)$ between $u_s(y+e_i)$ and $u_s(y)$ and square with factor
two. For the labelled source energy this proves
\begin{equation}\label{eq:graphcomparison}
 \En_p(u_s):=\sum_{y,i=1}^p|u_s(y+e_i)-u_s(y)|^2
 \leq2p\sum_y|u_s-v_s|^2+2\sum_{y,i}|u_s(y)-v_s(y+e_i)|^2
 \leq2pE_s.
\end{equation}
The repeated-rung factor is exactly $2p$, also at period two. Summing $HE_s$
over all layer pairs, both channels and all spatial components gives a lower
part of the original $Q_B$; weak-source terms are only discarded as positive.

\subsection{Balanced quantiles and both layer norms}
For an integer $1\leq m\leq n$, partition each source coordinate into
\[
 I_j=\{\lfloor jn/m\rfloor,\ldots,\lfloor(j+1)n/m\rfloor-1\},
 \qquad0\leq j<m.
\]
These are nonempty paths with lengths $\lfloor n/m\rfloor$ or
$L=\lceil n/m\rceil$. If $L=1$ all are singletons; otherwise a shorter side
is $L-1\geq L/2$. Every product cell therefore meets the exact hypotheses
of Lemma~\ref{lem:point}. Anchor its lowest-coordinate corner as a source
value $u_s(b)=0$, at the actual original point $x_s(b)$.
Cell internal edges are a subset of \eqref{eq:graphcomparison}, counted once;
across-cell and wrap edges are discarded as nonnegative.
Define
\[
 X_2=L^2(1+\ln L),\qquad X_3=L^3.
\]
Lemma~\ref{lem:point} and \eqref{eq:graphcomparison} give
$\norm{u_s}^2\leq52X_2E_s$ for $p=2$ and
$\norm{u_s}^2\leq678X_3E_s$ for $p=3$.
Rungs also give $\norm{v_s}^2\leq2\norm{u_s}^2+2\sum|v_s-u_s|^2$.
As $X_p\geq1$, both original layers satisfy
\begin{equation}\label{eq:pairnorm}
 \norm{u_s}^2+\norm{v_s}^2\leq
 \begin{cases}158X_2E_s,&p=2,\\2036X_3E_s,&p=3.\end{cases}
\end{equation}
Here $3\cdot52+2=158$ and $3\cdot678+2=2036$; no target-layer mass is omitted.
These formulas remain valid at singleton cells because the anchored source
value is zero and the rung still controls its target.

There are $(n/2)m^p$ distinct source anchors per channel: the cell corners
are distinct, \eqref{eq:paircoordinates} is a bijection, and different source
layers are disjoint. Both channels together have $A=nm^p$ channel-labelled
anchors. Use the \emph{same} anchor sites for all $D$ spatial components in
that channel. Although source-coordinate cells need not be original rectangles,
only their original \emph{point} sites are observed. Applying
Lemma~\ref{lem:isolation} independently to the two channels gives actual
canonical partitions with
\begin{equation}\label{eq:pointcount}
 N_{\rm out}\leq2D+2D^2A=2D+2D^2nm^p.
\end{equation}
The initial rows and every channel/component factor are included. Original
chart representatives, not wrapping rectangles or diagonal averages, realize
the points. A block-kernel field is zero at every required anchor; extra
nonpoint means only narrow this kernel.

\subsection{Dimension three: the critical logarithmic mesh}
Take $D=3$, $T_n=1+\ln n$, and let $m$ be the least positive integer with
$m^2\geq4096T_n$. If $m\leq n$, then $L\leq2n/m$ and $L\leq n$, giving
\[
 X_2=L^2(1+\ln L)\leq\frac{4n^2}{m^2}T_n\leq n^2/1024.
\]
Sum \eqref{eq:pairnorm} over the disjoint pairs, channels and components.
Since $H=cn^2$, a field zero at the prescribed anchors satisfies
\begin{equation}\label{eq:D3gap}
 Q_B(z)\geq\frac{H}{158X_2}\norm z^2
          \geq\frac{1024}{158}c\norm z^2\geq4c\norm z^2,
\end{equation}
because $4\cdot158=632<1024$. Put $\rho=\sqrt{4096T_n}\geq1$.
Minimality yields $m=\lceil\rho\rceil\leq2\rho$, so $m^2\leq16384T_n$.
The complete count \eqref{eq:pointcount} is therefore
\[
 N_{\rm out}\leq6+18nm^2\leq6+294912nT_n<2^{19}nT_n.
\]
We used $nT_n\geq1$ and $294912+6<524288$, without divisibility or exact-square
assumptions.

If $m>n$, minimality gives $(m-1)^2<4096T_n$ and the integers satisfy
$n\leq m-1$. The all-channel singleton fallback has
\[
 N_{\rm out}=6n^3<24576nT_n<2^{19}nT_n,\qquad K=c.
\]
This covers every ineligible even $n$, including $n=2$ and $n=m-1$.
At $m=n$ the mesh branch is allowed and its source cells are singleton;
\eqref{eq:D3gap} still follows from the derived rung bound. Thus no guessed
large-size cutoff, empty mesh cell or unproved singleton gap is used.

\subsection{Dimension four: a three-dimensional source mesh}
For $D=4$, first take $n\geq256$ and let $m$ be the least positive integer
with $m^3\geq65536n$. The integer $n$ itself is eligible because
$n^2\geq65536$, so $m\leq n$. Also $L\leq2n/m$ and
\[
 X_3=L^3\leq8n^3/m^3\leq n^2/8192.
\]
The other case of \eqref{eq:pairnorm} now gives, on the actual anchor kernel,
\begin{equation}\label{eq:D4gap}
 Q_B(z)\geq\frac{8192}{2036}c\norm z^2\geq4c\norm z^2,
 \qquad4\cdot2036=8144<8192.
\end{equation}
For $\rho=(65536n)^{1/3}\geq256$, $m=\lceil\rho\rceil\leq2\rho$, hence
$m^3\leq524288n$. Thus $A=nm^3\leq524288n^2$ and
\[
 N_{\rm out}\leq8+32A\leq8+2^{24}n^2\leq2^{25}n^2.
\]
At $n=256$ one has $m=n$ exactly; singleton source cells and their rungs
are valid in \eqref{eq:D4gap}. Noncubes are covered by the same ceiling
argument. For even $n<256$, use all-channel singletons:
$N_{\rm out}=8n^4<524288n^2<2^{25}n^2$ and $K=c$.
No three-dimensional source-cell estimate is applied in that fallback.

\subsection{Finite completion, constants and actual restricted rank}
For either mesh the canonical block kernel obeys
$Q_B(u)\geq4c\norm u^2$ by \eqref{eq:D3gap} or \eqref{eq:D4gap}.
Lemma~\ref{lem:completion} with $\delta=3c$ and $M=4DH+c$ proves the
all-field inequality with exactly \eqref{eq:pointK}. Dividing by $V$ and
using \eqref{eq:RC} gives the stated canonical normalization. On a constant
field $Q_B=0$ and $\mathcal Rz=z$; the finite coefficient is at least $c$,
so these constants are paid. The conservative mesh coefficient can grow as
$cn^4$; in the fallbacks $\mathcal R=I,K=c$ directly.
Restrict only after this all-field proof to the original $\UU_0^2$.
Its actual rank is at most the full count, with no assumed row independence
or coexact invariance of block averaging. This completes
Theorem~\ref{thm:checker}(ii)--(iii).

The same positive stack still has $D^2$ weak and $D^2$ strong eigenvalues per
site. Formula~\eqref{eq:fixedcap} with $q=1,g=2$, including its original
$\alpha_3=729/128$ or $\alpha_4=81/16$, is unchanged and has order $V$.
These improved observations do not reduce the inherited capped budget.
Neither the logarithm in this particular finite construction nor its point
count is asserted necessary for every permitted partition. The present
point proof is for \emph{orthogonal} lines only; it is distinct from the
fixed-angle plane result and supplies no claimed solution to its unproved
point-anchor analogue.
